% !TeX program = xelatex
% 由扫描件人工逐行转录。每个 \BookLine 对应原书一行，可直接编辑。
\documentclass[10pt]{article}
\usepackage[paperwidth=140mm,paperheight=210mm,margin=16mm]{geometry}
\usepackage{fontspec,xeCJK,graphicx,xcolor,tikz}
\usepackage[absolute,overlay]{textpos}
\usepackage[unicode,pdfpagelayout=TwoPageRight,pdftitle={阿尔刻提斯之歌（扫描部分重排）},pdfauthor={}]{hyperref}
\setmainfont{Liberation Serif}
\IfFileExists{/usr/local/share/fonts/f/FZShuSong_Z01_Regular.ttf}{
  \setCJKmainfont{FZShuSong_Z01_Regular.ttf}[Path=/usr/local/share/fonts/f/]
}{\setCJKmainfont{Noto Serif CJK SC}}
\IfFileExists{/usr/local/share/fonts/f/FZKai_Z03_Regular.ttf}{
  \newCJKfontfamily\kaifont{FZKai_Z03_Regular.ttf}[Path=/usr/local/share/fonts/f/]
}{\newCJKfontfamily\kaifont{FangSong}}
\IfFileExists{/usr/local/share/fonts/f/FZHei_B01_Regular.ttf}{
  \newCJKfontfamily\heifont{FZHei_B01_Regular.ttf}[Path=/usr/local/share/fonts/f/]
}{\newCJKfontfamily\heifont{Noto Sans CJK SC}}
\pagestyle{empty}
\setlength{\parindent}{0pt}
\setlength{\TPHorizModule}{1mm}
\setlength{\TPVertModule}{1mm}
\textblockorigin{0mm}{0mm}
\xeCJKsetup{PunctStyle=plain}
\newsavebox{\rowbox}
% y（毫米）, 缩进（毫米）, 字体, 是否铺满行, 文字。
\newcommand{\BookLine}[5]{%
  \begin{textblock*}{108mm}(16mm,#1mm)
  \fontsize{10.5}{14}\selectfont #3%
  \hspace*{#2mm}%
  \begingroup
  \ifnum#4=1 \renewcommand\CJKglue{\hskip 0pt plus 1fill}\fi
  \sbox{\rowbox}{#5}%
  \ifdim\wd\rowbox>\dimexpr108mm-#2mm\relax
    \resizebox{\dimexpr108mm-#2mm\relax}{!}{\usebox{\rowbox}}%
  \else
    \ifnum#4=1
      \makebox[\dimexpr108mm-#2mm\relax][s]{#5}%
    \else\usebox{\rowbox}\fi
  \fi
  \endgroup
  \end{textblock*}%
}
\newcommand{\Vertical}[5]{%
  \begin{textblock*}{8mm}(#1mm,#2mm)
    #3\fontsize{#4}{#5}\selectfont\centering
}
\newcommand{\EndVertical}{\end{textblock*}}
\newcommand{\BookNumber}[1]{%
  \begin{textblock*}{140mm}(0mm,192mm)
  \centering\fontsize{12}{12}\selectfont\itshape #1
  \end{textblock*}
}
\newcommand{\ActOrnament}[1][一]{%
  \Vertical{66}{28}{}{14}{21}第\\#1\\幕\EndVertical
  \begin{textblock*}{140mm}(0mm,63mm)
    \centering\textcolor{black!65}{\fontsize{10}{10}\selectfont\textbullet}
  \end{textblock*}
  \begin{textblock*}{75mm}(38mm,96mm)
  \begin{tikzpicture}[x=1mm,y=1mm,line width=.35pt,draw=black!65]
    % 后方拱门。
    \draw (16,0) -- (57,0) -- (57,39)
      .. controls (57,48) and (53,55) .. (48,59);
    \draw (16,0) -- (16,5);
    % 前方拱门。
    \draw (0,5) -- (48,5) -- (48,49)
      arc[start angle=0,end angle=180,radius=24] -- cycle;
    % 沿拱门上方逐字排出书名，不依赖扫描图。
    \foreach \letter [count=\i from 0] in {T,h,e,A,l,c,e,s,t,i,a,d}{%
      \pgfmathsetmacro{\a}{153-\i*11.45}
      \node[rotate=\a-90,text=black!80,font=\fontsize{9}{9}\selectfont]
        at ({24+35*cos(\a)},{49+35*sin(\a)}) {\letter};
    }
  \end{tikzpicture}
  \end{textblock*}
}
\newcommand{\WorkOrnament}[1]{%
  \Vertical{58}{26}{}{14}{21}#1\EndVertical
  \begin{textblock*}{140mm}(0mm,63mm)
    \centering\textcolor{black!65}{\fontsize{10}{10}\selectfont\textbullet}
  \end{textblock*}
  \begin{textblock*}{75mm}(38mm,96mm)
  \begin{tikzpicture}[x=1mm,y=1mm,line width=.35pt,draw=black!65]
    \draw (16,0) -- (57,0) -- (57,39) .. controls (57,48) and (53,55) .. (48,59);
    \draw (16,0) -- (16,5);
    \draw (0,5) -- (48,5) -- (48,49) arc[start angle=0,end angle=180,radius=24] -- cycle;
    \foreach \letter [count=\i from 0] in {T,h,e,A,l,c,e,s,t,i,a,d}{%
      \pgfmathsetmacro{\a}{153-\i*11.45}
      \node[rotate=\a-90,text=black!80,font=\fontsize{9}{9}\selectfont]
        at ({24+35*cos(\a)},{49+35*sin(\a)}) {\letter};
    }
  \end{tikzpicture}
  \end{textblock*}
}
\begin{document}
\null
\ActOrnament[二]
\newpage
% 原书页码：82
\null
\newpage
% 原书页码：83
\null
\BookNumber{83}
\Vertical{129}{11}{}{7}{10}第\\二\\幕\EndVertical
\BookLine{30.00}{7.38}{\kaifont}{0}{同一场景，十二年之后。仍然是黎明第一道光，仍然是}
\BookLine{38.50}{0.00}{\kaifont}{0}{同一个看门人绕过宫殿走出来。他并没有敲鼓，轻声而缓慢}
\BookLine{47.00}{0.00}{\kaifont}{0}{地说着话，情绪十分低落。}
\BookLine{64.00}{0.00}{\heifont}{0}{看门人：}
\BookLine{72.50}{22.14}{}{0}{黎明时分，照例巡查。热情诚挚的阿德墨托斯，}
\BookLine{81.00}{22.14}{}{0}{塞萨利之王，千金万马的拥有者。{\kaifont （低着头原地}}
\BookLine{89.50}{22.14}{}{0}{{\kaifont 站了一会儿，开始对观众说话）}事情要多糟糕}
\BookLine{98.00}{22.14}{}{0}{有多糟糕，非常糟糕。{\kaifont （门半掩着，他走到右}}
\BookLine{106.50}{22.14}{}{0}{{\kaifont 边门口，对外面的人们轻声说着）}没有，朋友}
\BookLine{115.00}{22.14}{}{0}{们，没什么新消息。国王还活着，他又撑过了}
\BookLine{123.50}{22.14}{}{0}{一晚，王后一直守着他，紧紧握着他的手。没}
\BookLine{132.00}{22.14}{}{0}{有新消息，没什么变化……把那条狗牵走，我}
\BookLine{140.50}{22.14}{}{0}{们这儿需要安静。对，国王喝了点掺着公牛血}
\BookLine{149.00}{22.14}{}{0}{的酒，{\kaifont （走回到场中央，对观众说道）}但是不得}
\BookLine{157.50}{22.14}{}{0}{不承认，国王已经处在死亡边缘。你们应该记}
\BookLine{166.00}{22.14}{}{0}{得那几个牧人，十二年前，忒瑞西阿斯带他们}
\BookLine{174.50}{22.14}{}{0}{到这儿来，说他们中间有一个人是阿波罗——}
\newpage
% 原书页码：84
\null
\BookNumber{84}
\Vertical{4}{169}{}{6.2}{8}阿\\尔\\刻\\提\\斯\\之\\歌\EndVertical
\BookLine{30.00}{22.14}{}{0}{阿波罗会作为仆人在这儿生活一年，你们还记}
\BookLine{38.50}{22.14}{}{0}{得吗？}
\BookLine{55.50}{22.14}{}{0}{嗯，一年结束之后他们走了吗？没有，他们四}
\BookLine{64.00}{22.14}{}{0}{个还待在这儿。阿波罗确实是这四人之一？没}
\BookLine{72.50}{22.14}{}{0}{人知道。别想这些事了，这种事想多了会变成}
\BookLine{81.00}{22.14}{}{0}{神经病的。不过朋友们，那一年确实很神奇。}
\BookLine{89.50}{22.14}{}{0}{我没法解释清楚：对于国家繁荣程度来说，那}
\BookLine{98.00}{22.14}{}{0}{一年和其他年份没太大区别。对于来访的，路}
\BookLine{106.50}{22.14}{}{0}{过的人来说，一切看起来都是老样子。但对我}
\BookLine{115.00}{22.14}{}{0}{们这些生活在其中的人来说，一切都不一样了。}
\BookLine{123.50}{22.14}{}{0}{人类生活的本质从没改变过，是我们看待生活}
\BookLine{132.00}{22.14}{}{0}{的眼光发生了改变。阿波罗确实就在这里。{\kaifont （停}}
\BookLine{140.50}{22.14}{}{0}{{\kaifont 顿）}嗯？阿德墨托斯国王怎么生病了？我不是}
\BookLine{149.00}{22.14}{}{0}{说了吗，那四个牧人留了下来。他们是好人，}
\BookLine{157.50}{22.14}{}{0}{好工人——虽然他们操着一口希腊南方的乡村}
\BookLine{166.00}{22.14}{}{0}{口音。其中还有两人——呃！嗜酒如命。有一}
\BookLine{174.50}{22.14}{}{0}{天晚上，太阳落山的时候他们就坐在那儿喝酒，}
\newpage
% 原书页码：85
\null
\BookNumber{85}
\Vertical{129}{11}{}{7}{10}第\\二\\幕\EndVertical
\BookLine{30.00}{22.14}{}{0}{那儿，就在那儿，一边喝一边开始吵架。他们}
\BookLine{38.50}{22.14}{}{0}{叫着嚷着，阿德墨托斯国王出来想要制止。结}
\BookLine{47.00}{22.14}{}{0}{果，天哪，其中一个牧人拿着刀，一不小心刺}
\BookLine{55.50}{22.14}{}{0}{到了国王，从这儿{\kaifont （指着喉咙）}一刀划到腰上。}
\BookLine{64.00}{22.14}{}{0}{伤口大得能把手伸进去。这已经是几星期前的}
\BookLine{72.50}{22.14}{}{0}{事儿了，但伤口还是没见好。伤口发炎然后结}
\BookLine{81.00}{22.14}{}{0}{痂，又流水又结痂的，国王看来是难逃此劫了。}
\BookLine{89.50}{22.14}{}{0}{没人打算怪罪那个牧人，因为他可能就是阿波}
\BookLine{98.00}{22.14}{}{0}{罗啊——你明白吗？这或许就是阿波罗的旨}
\BookLine{106.50}{22.14}{}{0}{意——想想看！现在那个牧人就在门外——他}
\BookLine{115.00}{22.14}{}{0}{一直在那儿——跪在地上，把脸埋在地上，希}
\BookLine{123.50}{22.14}{}{0}{望死的是他自己。}
\BookLine{140.50}{7.38}{\kaifont}{0}{门边有人在嘀咕，声音越来越大，听上去很激动。}
\BookLine{157.50}{0.00}{\heifont}{0}{看门人：}
\BookLine{166.00}{22.14}{}{0}{这又是怎么了？}
\newpage
% 原书页码：86
\null
\BookNumber{86}
\Vertical{4}{169}{}{6.2}{8}阿\\尔\\刻\\提\\斯\\之\\歌\EndVertical
\BookLine{30.00}{7.38}{\kaifont}{0}{阿格莱亚急匆匆地从宫殿里出来。}
\BookLine{47.00}{0.00}{\heifont}{0}{阿格莱亚：}
\BookLine{55.50}{22.14}{}{0}{出什么事了？看门人，这么吵闹你也不管管？}
\BookLine{72.50}{0.00}{\heifont}{0}{看门人：}
\BookLine{81.00}{22.14}{}{0}{{\kaifont （一瘸一拐地朝门口走去）}我也刚听见。}
\BookLine{98.00}{0.00}{\heifont}{0}{阿格莱亚：}
\BookLine{106.50}{22.14}{}{0}{{\kaifont （站在门边）}你们啰唆个不停在说什么？不知道}
\BookLine{115.00}{22.14}{}{0}{我们这儿需要安静吗？}
\BookLine{132.00}{0.00}{\heifont}{0}{看门人：}
\BookLine{140.50}{22.14}{}{0}{你们讲不讲理啊？你们有没有心肝啊？什么？}
\BookLine{149.00}{22.14}{}{0}{什么？}
\BookLine{166.00}{0.00}{\heifont}{0}{阿格莱亚：}
\BookLine{174.50}{22.14}{}{0}{什么？什么信使？告诉他去宫殿前门，一定要}
\newpage
% 原书页码：87
\null
\BookNumber{87}
\Vertical{129}{11}{}{7}{10}第\\二\\幕\EndVertical
\BookLine{30.00}{22.14}{}{0}{安静。什么信使？从哪儿来的？}
\BookLine{47.00}{7.38}{\kaifont}{0}{阿尔刻提斯上场，她正从宫殿里走出来。}
\BookLine{64.00}{0.00}{\heifont}{0}{阿尔刻提斯：}
\BookLine{72.50}{22.14}{}{0}{我们这儿要安静，要安静。阿格莱亚，这么吵}
\BookLine{81.00}{22.14}{}{0}{闹你也不管管吗？}
\BookLine{98.00}{0.00}{\heifont}{0}{阿格莱亚：}
\BookLine{106.50}{22.14}{}{0}{阿尔刻提斯王后，他们说来了一位信使。}
\BookLine{123.50}{0.00}{\heifont}{0}{阿尔刻提斯：}
\BookLine{132.00}{22.14}{}{0}{什么信使？}
\BookLine{149.00}{7.38}{\kaifont}{0}{有人从门缝递给他一小片长方形的金叶子，他转交给阿}
\BookLine{157.50}{0.00}{\kaifont}{0}{尔刻提斯。}
\newpage
% 原书页码：88
\null
\BookNumber{88}
\Vertical{4}{169}{}{6.2}{8}阿\\尔\\刻\\提\\斯\\之\\歌\EndVertical
\BookLine{30.00}{0.00}{\heifont}{0}{看门人：}
\BookLine{38.50}{22.14}{}{0}{他们说夜里来了一位信使，留下了这个。}
\BookLine{55.50}{7.38}{\kaifont}{0}{阿尔刻提斯拿在手里看了看，看门人也凑上去仔细看了}
\BookLine{64.00}{0.00}{\kaifont}{0}{看，严肃地说道。}
\BookLine{81.00}{0.00}{\heifont}{0}{看门人：}
\BookLine{89.50}{22.14}{}{0}{这是金的，夫人。您看，那些标志——是南方}
\BookLine{98.00}{22.14}{}{0}{的文字。那是太阳——祭坛，还有桂冠……阿}
\BookLine{106.50}{22.14}{}{0}{尔刻提斯王后，这是从德尔斐送来的，这是从}
\BookLine{115.00}{22.14}{}{0}{德尔斐送来的神谕！}
\BookLine{132.00}{7.38}{\kaifont}{0}{阿尔刻提斯立刻行起了仪式：她把叶子放到额前，然后}
\BookLine{140.50}{0.00}{\kaifont}{0}{放到胸前，然后放到唇边。}
\BookLine{157.50}{0.00}{\heifont}{0}{阿尔刻提斯：}
\BookLine{166.00}{22.14}{}{0}{信使在哪里？}
\newpage
% 原书页码：89
\null
\BookNumber{89}
\Vertical{129}{11}{}{7}{10}第\\二\\幕\EndVertical
\BookLine{30.00}{7.38}{\kaifont}{0}{看门人走到此刻微微开着的门边，小声与外面的人交谈}
\BookLine{38.50}{0.00}{\kaifont}{0}{起来。}
\BookLine{55.50}{0.00}{\heifont}{0}{看门人：}
\BookLine{64.00}{22.14}{}{0}{{\kaifont （转身对阿尔刻提斯）}他们说，他几个小时之前}
\BookLine{72.50}{22.14}{}{0}{走了。}
\BookLine{89.50}{0.00}{\heifont}{0}{阿尔刻提斯：}
\BookLine{98.00}{22.14}{}{0}{谁能看懂这些字？看门人，去把那四个牧人}
\BookLine{106.50}{22.14}{}{0}{叫来！}
\BookLine{123.50}{0.00}{\heifont}{0}{看门人：}
\BookLine{132.00}{22.14}{}{0}{他们就在这儿，夫人，四个人都在这儿了。}
\BookLine{149.00}{7.38}{\kaifont}{0}{门口又传来说话的声音。看门人开门让四个牧人进来。}
\BookLine{157.50}{0.00}{\kaifont}{0}{第一幕曾与阿尔刻提斯长谈过的那个牧人俯伏在地上，脸朝}
\BookLine{166.00}{0.00}{\kaifont}{0}{着地面，双手抱着自己的头。}
\newpage
% 原书页码：90
\null
\BookNumber{90}
\Vertical{4}{169}{}{6.2}{8}阿\\尔\\刻\\提\\斯\\之\\歌\EndVertical
\BookLine{30.00}{0.00}{\heifont}{0}{阿尔刻提斯：}
\BookLine{38.50}{22.14}{}{0}{你们有人认识南方文字吗？}
\BookLine{55.50}{0.00}{\heifont}{0}{牧人：}
\BookLine{64.00}{22.14}{}{0}{{\kaifont （抬起头）}我认识，一点，阿尔刻提斯王后。}
\BookLine{81.00}{0.00}{\heifont}{0}{阿尔刻提斯：}
\BookLine{89.50}{22.14}{}{0}{你先起来。站起来。}
\BookLine{106.50}{7.38}{\kaifont}{0}{他站了起来，她把叶子放在他手里，他看了一下，然后}
\BookLine{115.00}{0.00}{\kaifont}{0}{把叶子依次放在自己额前、胸前和唇边，敬畏地说道。}
\BookLine{132.00}{0.00}{\heifont}{0}{牧人：}
\BookLine{140.50}{22.14}{}{0}{这是从德尔斐的阿波罗神殿送来的信息。}
\BookLine{157.50}{7.38}{\kaifont}{0}{几个镇上的人从门口挤了进来，牧人艰难地辨认并读出}
\BookLine{166.00}{0.00}{\kaifont}{0}{文字。}
\newpage
% 原书页码：91
\null
\BookNumber{91}
\Vertical{129}{11}{}{7}{10}第\\二\\幕\EndVertical
\BookLine{30.00}{0.00}{\heifont}{0}{牧人：}
\BookLine{38.50}{22.14}{}{0}{平安……与长寿……归于热情诚挚的阿德墨托}
\BookLine{47.00}{22.14}{}{0}{斯，塞萨利之王，千金万马的拥有者。}
\BookLine{64.00}{0.00}{\heifont}{0}{阿尔刻提斯：}
\BookLine{72.50}{22.14}{}{0}{长寿？}
\BookLine{89.50}{0.00}{\heifont}{0}{牧人：}
\BookLine{98.00}{22.14}{}{0}{平安与长寿归于热情诚挚的阿德墨托斯。}
\BookLine{115.00}{0.00}{\heifont}{0}{阿尔刻提斯：}
\BookLine{123.50}{22.14}{}{0}{等等。{\kaifont （对看门人）}把门关起来。}
\BookLine{140.50}{7.38}{\kaifont}{0}{看门人把镇上的人赶下台去，他看了阿尔刻提斯一眼，}
\BookLine{149.00}{0.00}{\kaifont}{0}{把其他三个牧人也叫了下去，然后关上门。阿尔刻提斯走到}
\BookLine{157.50}{0.00}{\kaifont}{0}{牧人身旁，他正在努力辨认叶子上的字。}
\newpage
% 原书页码：92
\null
\BookNumber{92}
\Vertical{4}{169}{}{6.2}{8}阿\\尔\\刻\\提\\斯\\之\\歌\EndVertical
\BookLine{30.00}{0.00}{\heifont}{0}{牧人：}
\BookLine{38.50}{22.14}{}{0}{“阿德墨托斯国王，不，不会，死。”}
\BookLine{55.50}{0.00}{\heifont}{0}{阿格莱亚：}
\BookLine{64.00}{22.14}{}{0}{伟大的阿波罗！伟大的阿波罗！}
\BookLine{81.00}{0.00}{\heifont}{0}{牧人：}
\BookLine{89.50}{22.14}{}{0}{“不会死，如果，如果，因为，不，如果，另一}
\BookLine{98.00}{22.14}{}{0}{个，如果另一个人，第二个人，想要，{\kaifont （他指了}}
\BookLine{106.50}{22.14}{}{0}{{\kaifont 指自己的心）}希望，渴望，想要去死。”夫人，}
\BookLine{115.00}{22.14}{}{0}{我不认识最后一个字。}
\BookLine{132.00}{0.00}{\heifont}{0}{阿尔刻提斯：}
\BookLine{140.50}{22.14}{}{0}{我明白了。}
\BookLine{157.50}{0.00}{\heifont}{0}{牧人：}
\BookLine{166.00}{22.14}{}{0}{{\kaifont （费力地，终于想到）}“在他的位置，代替他。”}
\BookLine{174.50}{22.14}{}{0}{伟大的阿波罗！}
\newpage
% 原书页码：93
\null
\BookNumber{93}
\Vertical{129}{11}{}{7}{10}第\\二\\幕\EndVertical
\BookLine{30.00}{0.00}{\heifont}{0}{阿格莱亚：}
\BookLine{38.50}{22.14}{}{0}{这是什么意思？什么意思？}
\BookLine{55.50}{0.00}{\heifont}{0}{看门人：}
\BookLine{64.00}{22.14}{}{0}{我。我应该去死。{\kaifont （走开）}我的剑，我的剑呢？}
\BookLine{81.00}{0.00}{\heifont}{0}{阿格莱亚：}
\BookLine{89.50}{22.14}{}{1}{{\kaifont （明白过来）}不。{\kaifont （对看门人坚决地说道）}不，这}
\BookLine{98.00}{22.14}{}{0}{是给我的信息。我是看着他出生的，我也应该}
\BookLine{106.50}{22.14}{}{0}{为了他去死。我要去投河自尽。}
\BookLine{123.50}{0.00}{\heifont}{0}{看门人：}
\BookLine{132.00}{22.14}{}{0}{{\kaifont （回答阿格莱亚）}这不是女人该干的事情。}
\BookLine{149.00}{0.00}{\heifont}{0}{牧人：}
\BookLine{157.50}{22.14}{}{0}{{\kaifont （跪下，他用双拳捂住自己的眼睛）}我早就想要}
\BookLine{166.00}{22.14}{}{0}{去死了。}
\newpage
% 原书页码：94
\null
\BookNumber{94}
\Vertical{4}{169}{}{6.2}{8}阿\\尔\\刻\\提\\斯\\之\\歌\EndVertical
\BookLine{30.00}{0.00}{\heifont}{0}{阿格莱亚：}
\BookLine{38.50}{22.14}{}{0}{阿尔刻提斯王后，请吩咐吧，我一定办到。要}
\BookLine{47.00}{22.14}{}{0}{怎么为阿德墨托斯国王去死？}
\BookLine{64.00}{0.00}{\heifont}{0}{看门人：}
\BookLine{72.50}{22.14}{}{0}{{\kaifont （自告奋勇，站在阿格莱亚和阿尔刻提斯中间）}}
\BookLine{81.00}{22.14}{}{0}{阿尔刻提斯王后，下命令给我吧。很明显这个}
\BookLine{89.50}{22.14}{}{0}{信息就是给我的。}
\BookLine{106.50}{0.00}{\heifont}{0}{牧人：}
\BookLine{115.00}{22.14}{}{0}{{\kaifont （再次全身俯伏在地）}是我伤害了他，这是我应}
\BookLine{123.50}{22.14}{}{0}{得的！}
\BookLine{140.50}{0.00}{\heifont}{0}{阿尔刻提斯：}
\BookLine{149.00}{22.14}{}{0}{{\kaifont （一直一动不动，沉默着没有说话）}你们怎么能}
\BookLine{157.50}{22.14}{}{0}{替他去死呢？{\kaifont （停顿）}你们以为把自己扔进河里}
\BookLine{166.00}{22.14}{}{0}{就行了吗？一刀刺进心里就行了吗？不，神明}
\BookLine{174.50}{22.14}{}{0}{们不会让一切来得这么容易。}
\newpage
% 原书页码：95
\null
\BookNumber{95}
\Vertical{129}{11}{}{7}{10}第\\二\\幕\EndVertical
\BookLine{30.00}{0.00}{\heifont}{0}{阿格莱亚：}
\BookLine{38.50}{22.14}{}{0}{{\kaifont （啜泣着）}阿尔刻提斯王后，要怎么为阿德墨托}
\BookLine{47.00}{22.14}{}{0}{斯国王去死？}
\BookLine{64.00}{0.00}{\heifont}{0}{牧人：}
\BookLine{72.50}{22.14}{}{0}{我知道，我知道。{\kaifont （转身向门口走去）}平安与长}
\BookLine{81.00}{22.14}{}{0}{寿归于阿德墨托斯国王！}
\BookLine{98.00}{0.00}{\heifont}{0}{阿尔刻提斯：}
\BookLine{106.50}{22.14}{}{0}{{\kaifont （对牧人）}站住！}
\BookLine{123.50}{7.38}{\kaifont}{0}{大家都看着她，她平静地说道。}
\BookLine{140.50}{0.00}{\heifont}{0}{阿尔刻提斯：}
\BookLine{149.00}{22.14}{}{0}{阿格莱亚，去把我结婚时穿的那条裙子拿出来。}
\BookLine{166.00}{0.00}{\heifont}{0}{阿格莱亚：}
\BookLine{174.50}{22.14}{}{0}{{\kaifont （突然明白过来，盯着她，惊恐地小声说）}不，}
\newpage
% 原书页码：96
\null
\BookNumber{96}
\Vertical{4}{169}{}{6.2}{8}阿\\尔\\刻\\提\\斯\\之\\歌\EndVertical
\BookLine{30.00}{22.14}{}{0}{不，您不能……}
\BookLine{47.00}{0.00}{\heifont}{0}{阿尔刻提斯：}
\BookLine{55.50}{22.14}{}{0}{告诉孩子们我等下过来看他们。告诉国王外面}
\BookLine{64.00}{22.14}{}{0}{太阳很暖和，他应该出来晒晒太阳。}
\BookLine{81.00}{7.38}{\kaifont}{0}{大家彼此争论起来。}
\BookLine{98.00}{0.00}{\heifont}{0}{阿格莱亚：}
\BookLine{106.50}{22.14}{}{0}{不，不，阿尔刻提斯王后。我们已经老了，我}
\BookLine{115.00}{22.14}{}{0}{们活得差不多了。}
\BookLine{132.00}{0.00}{\heifont}{0}{看门人：}
\BookLine{140.50}{22.14}{}{0}{夫人，您看看我——我是个老人。}
\BookLine{157.50}{0.00}{\heifont}{0}{阿格莱亚：}
\BookLine{166.00}{22.14}{}{0}{您不能这么做，他也不会同意的。您是王后，}
\BookLine{174.50}{22.14}{}{0}{您是孩子们的母亲。}
\newpage
% 原书页码：97
\null
\BookNumber{97}
\Vertical{129}{11}{}{7}{10}第\\二\\幕\EndVertical
\BookLine{30.00}{0.00}{\heifont}{0}{阿尔刻提斯：}
\BookLine{38.50}{22.14}{}{0}{{\kaifont （同时用言语和手势制止他们的发言）}首先！首}
\BookLine{47.00}{22.14}{}{0}{先我命令你们，这件事情不许对任何人说起。}
\BookLine{55.50}{22.14}{}{0}{阿格莱亚，你听到了吗？}
\BookLine{72.50}{0.00}{\heifont}{0}{阿格莱亚：}
\BookLine{81.00}{22.14}{}{0}{是的，阿尔刻提斯王后。}
\BookLine{98.00}{0.00}{\heifont}{0}{阿尔刻提斯：}
\BookLine{106.50}{22.14}{}{0}{看门人！}
\BookLine{123.50}{0.00}{\heifont}{0}{看门人：}
\BookLine{132.00}{22.14}{}{0}{是的，阿尔刻提斯王后。}
\BookLine{149.00}{0.00}{\heifont}{0}{阿尔刻提斯：}
\BookLine{157.50}{22.14}{}{0}{直到明天，不能对任何人说。无论这里发生了}
\BookLine{166.00}{22.14}{}{0}{什么，你们都不许惊讶，不许悲伤。现在让我}
\BookLine{174.50}{22.14}{}{0}{和他单独待一会儿。{\kaifont （指着牧人）}}
\newpage
% 原书页码：98
\null
\BookNumber{98}
\Vertical{4}{169}{}{6.2}{8}阿\\尔\\刻\\提\\斯\\之\\歌\EndVertical
\BookLine{30.00}{7.38}{\kaifont}{0}{阿格莱亚走进宫殿，看门人下场。牧人站在那里，一只}
\BookLine{38.50}{0.00}{\kaifont}{0}{手放在额前沉思着，突然喊了起来。}
\BookLine{55.50}{0.00}{\heifont}{0}{牧人：}
\BookLine{64.00}{22.14}{}{0}{阿德墨托斯国王——起来吧！起来吧！}
\BookLine{81.00}{0.00}{\heifont}{0}{阿尔刻提斯：}
\BookLine{89.50}{22.14}{}{0}{你救不了他的。}
\BookLine{106.50}{0.00}{\heifont}{0}{牧人：}
\BookLine{115.00}{22.14}{}{0}{我早就想要去死了。}
\BookLine{132.00}{0.00}{\heifont}{0}{阿尔刻提斯：}
\BookLine{140.50}{22.14}{}{0}{是，或许你能做到，但你没法把这件事完全做}
\BookLine{149.00}{22.14}{}{0}{好。你希望去死，是的，但不是因为你爱阿德}
\BookLine{157.50}{22.14}{}{0}{墨托斯，只是为了让心中的罪恶感得到解脱。}
\BookLine{166.00}{22.14}{}{0}{这件事应该是为了爱而做，不是为了赎罪，是}
\BookLine{174.50}{22.14}{}{0}{为了爱。}
\newpage
% 原书页码：99
\null
\BookNumber{99}
\Vertical{129}{11}{}{7}{10}第\\二\\幕\EndVertical
\BookLine{30.00}{0.00}{\heifont}{0}{牧人：}
\BookLine{38.50}{22.14}{}{0}{{\kaifont （有些恼怒）}我，也爱着阿德墨托斯。}
\BookLine{55.50}{0.00}{\heifont}{0}{阿尔刻提斯：}
\BookLine{64.00}{22.14}{}{0}{有谁不爱他呢？但你的死太轻微了。你渴望去}
\BookLine{72.50}{22.14}{}{0}{死，而我，恐惧、害怕、憎恶死亡。我必须为}
\BookLine{81.00}{22.14}{}{0}{了阿德墨托斯而死，{\kaifont （望着天空）}在这日光下。}
\BookLine{89.50}{22.14}{}{0}{只有这样，他才能够重生。你能为他做到这}
\BookLine{98.00}{22.14}{}{0}{些吗？}
\BookLine{115.00}{7.38}{\kaifont}{0}{牧人沉默了，她继续自言自语般地说着。}
\BookLine{132.00}{22.14}{}{0}{我现在知道了要做什么，该怎么做，但我不明}
\BookLine{140.50}{22.14}{}{0}{白为什么……为什么是我。你……你能让我明}
\BookLine{149.00}{22.14}{}{0}{白，为什么我必须死去。}
\BookLine{166.00}{0.00}{\heifont}{0}{牧人：}
\BookLine{174.50}{22.14}{}{0}{我？}
\newpage
% 原书页码：100
\null
\BookNumber{100}
\Vertical{4}{169}{}{6.2}{8}阿\\尔\\刻\\提\\斯\\之\\歌\EndVertical
\BookLine{30.00}{0.00}{\heifont}{0}{阿尔刻提斯：}
\BookLine{38.50}{22.14}{}{0}{是的，你，你从德尔斐来，阿波罗就在你……}
\BookLine{55.50}{0.00}{\heifont}{0}{牧人：}
\BookLine{64.00}{22.14}{}{0}{{\kaifont （真诚地回绝）}夫人，公主，我告诉过您很多}
\BookLine{72.50}{22.14}{}{0}{很多次了。就算阿波罗在这儿，也绝对和我没}
\BookLine{81.00}{22.14}{}{0}{关系。}
\BookLine{98.00}{0.00}{\heifont}{0}{阿尔刻提斯：}
\BookLine{106.50}{22.14}{}{0}{{\kaifont （温和而轻柔地）}是你让我明白了，我的婚姻是}
\BookLine{115.00}{22.14}{}{0}{阿波罗的旨意。}
\BookLine{132.00}{0.00}{\heifont}{0}{牧人：}
\BookLine{140.50}{22.14}{}{0}{{\kaifont （有些恼怒）}不，不。您当时问了问题，我像其}
\BookLine{149.00}{22.14}{}{0}{他人一样回答您。您知道，我不过是个普普通}
\BookLine{157.50}{22.14}{}{0}{通的人。}
\BookLine{174.50}{7.38}{\kaifont}{0}{阿尔刻提斯转过身去，面对宫殿大门，沮丧至极，片刻}
\newpage
% 原书页码：101
\null
\BookNumber{101}
\Vertical{129}{11}{}{7}{10}第\\二\\幕\EndVertical
\BookLine{30.00}{0.00}{\kaifont}{0}{又转过身来对牧人说话。}
\BookLine{47.00}{0.00}{\heifont}{0}{阿尔刻提斯：}
\BookLine{55.50}{22.14}{}{0}{那就再次作为一个普普通通的人回答我的问题，}
\BookLine{64.00}{22.14}{}{0}{告诉我为什么阿波罗要让我死。}
\BookLine{81.00}{0.00}{\heifont}{0}{牧人：}
\BookLine{89.50}{22.14}{}{0}{{\kaifont （仍然带着一些恼怒）}那么作为一个普普通通的}
\BookLine{98.00}{22.14}{}{0}{人，我现在回答您——像其他人一样回答您：}
\BookLine{106.50}{22.14}{}{0}{德尔斐的旨意说，我们其中一人必须去死，我}
\BookLine{115.00}{22.14}{}{0}{也做好准备去死了，为什么一定要明白为什么}
\BookLine{123.50}{22.14}{}{0}{呢？}
\BookLine{140.50}{0.00}{\heifont}{0}{阿尔刻提斯：}
\BookLine{149.00}{22.14}{}{0}{如果不明白原因，那我们的生命和动物有什么}
\BookLine{157.50}{22.14}{}{0}{区别？}
\newpage
% 原书页码：102
\null
\BookNumber{102}
\Vertical{4}{169}{}{6.2}{8}阿\\尔\\刻\\提\\斯\\之\\歌\EndVertical
\BookLine{30.00}{0.00}{\heifont}{0}{牧人：}
\BookLine{38.50}{22.14}{}{0}{不！公主，明白，意味着您要看到整个事情。}
\BookLine{47.00}{22.14}{}{0}{我们人类有谁看到过整个故事呢，有谁看到过}
\BookLine{55.50}{22.14}{}{0}{故事的结局呢？如果您让我为阿德墨托斯去}
\BookLine{64.00}{22.14}{}{0}{死，我不知道等我死后会发生什么，但是我心}
\BookLine{72.50}{22.14}{}{0}{甘情愿。因为我一直明白这世上有两种死亡：}
\BookLine{81.00}{22.14}{}{0}{一种是终局；另一种是延续，因着随之而来的}
\BookLine{89.50}{22.14}{}{0}{一切而有了分量。我知道，如果这死亡的命运}
\BookLine{98.00}{22.14}{}{0}{是从德尔斐赐下给我的，那这死就不会是白白}
\BookLine{106.50}{22.14}{}{0}{的。如果众神真的存在，我想这就是他们的启}
\BookLine{115.00}{22.14}{}{0}{示：众神所做的一切还未曾揭示——而我们看}
\BookLine{123.50}{22.14}{}{0}{到的只是那广大之中极小的一部分。让我接受}
\BookLine{132.00}{22.14}{}{0}{这死亡吧，公主。因为这能够拯救我脱离另外}
\BookLine{140.50}{22.14}{}{0}{一种死亡，那正是我所恐惧的，也是所有人恐}
\BookLine{149.00}{22.14}{}{0}{惧的——死，不过是生命的停止，坟墓中的}
\BookLine{157.50}{22.14}{}{0}{尘埃。}
\BookLine{174.50}{7.38}{\kaifont}{0}{阿尔刻提斯在他的言辞中找到了所有答案。她的情绪有}
\newpage
% 原书页码：103
\null
\BookNumber{103}
\Vertical{129}{11}{}{7}{10}第\\二\\幕\EndVertical
\BookLine{30.00}{0.00}{\kaifont}{0}{了变化，轻声而迅速地说道。}
\BookLine{47.00}{0.00}{\heifont}{0}{阿尔刻提斯：}
\BookLine{55.50}{22.14}{}{0}{不，牧人，活下去——为了阿德墨托斯，为了}
\BookLine{64.00}{22.14}{}{0}{我，为了我的孩子们活下去。你是伊皮伊尼兹}
\BookLine{72.50}{22.14}{}{0}{的朋友呢，虽然因为你伤害了他父亲，他现在}
\BookLine{81.00}{22.14}{}{0}{恨你，你让他伤心坏了。他一直把你当作好朋}
\BookLine{89.50}{22.14}{}{0}{友，是你教会了他游泳、钓鱼。在离开之前，}
\BookLine{98.00}{22.14}{}{0}{我会告诉伊皮伊尼兹，你给了我多么大的恩惠。}
\BookLine{115.00}{7.38}{\kaifont}{0}{牧人默默走了出去，阿尔刻提斯拿金叶子又行了一次}
\BookLine{123.50}{0.00}{\kaifont}{0}{仪式：额前，胸前，唇边。阿格莱亚从宫殿里走了出来。}
\BookLine{140.50}{0.00}{\heifont}{0}{阿格莱亚：}
\BookLine{149.00}{22.14}{}{0}{{\kaifont （低着头）}裙子已经准备好了，阿尔刻提斯}
\BookLine{157.50}{22.14}{}{0}{王后。}
\newpage
% 原书页码：104
\null
\BookNumber{104}
\Vertical{4}{169}{}{6.2}{8}阿\\尔\\刻\\提\\斯\\之\\歌\EndVertical
\BookLine{30.00}{0.00}{\heifont}{0}{阿尔刻提斯：}
\BookLine{38.50}{22.14}{}{0}{阿格莱亚，我现在要做这件事，所以我不能，}
\BookLine{47.00}{22.14}{}{0}{绝不能去见孩子们。他们的小脑瓜闻起来太美}
\BookLine{55.50}{22.14}{}{0}{好了。你明白吗？}
\BookLine{72.50}{0.00}{\heifont}{0}{阿格莱亚：}
\BookLine{81.00}{22.14}{}{0}{明白，阿尔刻提斯王后。}
\BookLine{98.00}{0.00}{\heifont}{0}{阿尔刻提斯：}
\BookLine{106.50}{22.14}{}{0}{你告诉伊皮伊尼兹，就说我说的，让他去见那}
\BookLine{115.00}{22.14}{}{0}{个刺伤了他父亲的牧人，让他原谅他，并感谢}
\BookLine{123.50}{22.14}{}{0}{他给予我的恩惠。}
\BookLine{140.50}{0.00}{\heifont}{0}{阿格莱亚：}
\BookLine{149.00}{22.14}{}{0}{我会告诉他。}
\BookLine{166.00}{0.00}{\heifont}{0}{阿尔刻提斯：}
\BookLine{174.50}{22.14}{}{0}{你帮我剪下一束头发来，不要告诉任何人，把}
\newpage
% 原书页码：105
\null
\BookNumber{105}
\Vertical{129}{11}{}{7}{10}第\\二\\幕\EndVertical
\BookLine{30.00}{22.14}{}{0}{它放在祭坛上，{\kaifont （指着门）}就在路的那边，阿波}
\BookLine{38.50}{22.14}{}{0}{罗的丛林里。阿格莱亚，等我走了之后，告诉}
\BookLine{47.00}{22.14}{}{0}{阿德墨托斯国王，我希望他再娶一位妻子。}
\BookLine{64.00}{0.00}{\heifont}{0}{阿格莱亚：}
\BookLine{72.50}{22.14}{}{0}{王后啊！}
\BookLine{89.50}{0.00}{\heifont}{0}{阿尔刻提斯：}
\BookLine{98.00}{22.14}{}{0}{一个男人需要这样的安慰。但是，阿格莱亚！}
\BookLine{106.50}{22.14}{}{0}{人们都说继母会对前妻留下的孩子非常不好。}
\BookLine{115.00}{22.14}{}{0}{你要陪着他们，看好他们。阿格莱亚——要时}
\BookLine{123.50}{22.14}{}{0}{不时地，对他提起我。{\kaifont （快要说不下去了）}对他}
\BookLine{132.00}{22.14}{}{0}{提起我！}
\BookLine{149.00}{7.38}{\kaifont}{0}{阿尔刻提斯跑进宫殿里。阿格莱亚正要跟上去，听到门}
\BookLine{157.50}{0.00}{\kaifont}{0}{口有人在兴高采烈地讨论着什么。一些镇上的人走进了院子，}
\BookLine{166.00}{0.00}{\kaifont}{0}{罗德普，宫里的一个年轻侍女，一边掩饰着欢笑，一边从门}
\BookLine{174.50}{0.00}{\kaifont}{0}{口跑向宫殿大门。}
\newpage
% 原书页码：106
\null
\BookNumber{106}
\Vertical{4}{169}{}{6.2}{8}阿\\尔\\刻\\提\\斯\\之\\歌\EndVertical
\BookLine{30.00}{0.00}{\heifont}{0}{阿格莱亚：}
\BookLine{38.50}{22.14}{}{0}{有什么事吗？外面吵吵闹闹是怎么回事，罗}
\BookLine{47.00}{22.14}{}{0}{德普！}
\BookLine{64.00}{0.00}{\heifont}{0}{罗德普：}
\BookLine{72.50}{22.14}{}{0}{他马上就要到了！}
\BookLine{89.50}{0.00}{\heifont}{0}{阿格莱亚：}
\BookLine{98.00}{22.14}{}{0}{谁？}
\BookLine{115.00}{0.00}{\heifont}{0}{罗德普和镇上的人们：}
\BookLine{123.50}{22.14}{}{0}{赫拉克勒斯来了！他已经从山谷下来了！宙斯}
\BookLine{132.00}{22.14}{}{0}{之子，赫拉克勒斯！}
\BookLine{149.00}{0.00}{\heifont}{0}{阿格莱亚：}
\BookLine{157.50}{22.14}{}{0}{你们都闭嘴。快从院子里出去。你们难道忘了}
\BookLine{166.00}{22.14}{}{0}{这个家里有人生病吗？罗德普，快给我进去，}
\BookLine{174.50}{22.14}{}{0}{闭上你的嘴！}
\newpage
% 原书页码：107
\null
\BookNumber{107}
\Vertical{129}{11}{}{7}{10}第\\二\\幕\EndVertical
\BookLine{30.00}{7.38}{\kaifont}{0}{镇上的人们下场。看门人打开宫殿的大门，罗德普刚好}
\BookLine{38.50}{0.00}{\kaifont}{0}{溜了进去。看门人身后跟着两个仆人，被他们抬着的卧榻挡}
\BookLine{47.00}{0.00}{\kaifont}{0}{着。他对阿格莱亚说话，而她正焦虑地低着头。}
\BookLine{64.00}{0.00}{\heifont}{0}{看门人：}
\BookLine{72.50}{22.14}{}{0}{赫拉克勒斯！大日子！}
\BookLine{89.50}{0.00}{\heifont}{0}{阿格莱亚：}
\BookLine{98.00}{22.14}{}{0}{{\kaifont （讽刺着）}是啊，是啊，每年都来。他为什么每}
\BookLine{106.50}{22.14}{}{0}{年都要来，为什么？}
\BookLine{123.50}{0.00}{\heifont}{0}{看门人：}
\BookLine{132.00}{22.14}{}{0}{这有什么，多尊贵的朋友！}
\BookLine{149.00}{0.00}{\heifont}{0}{阿格莱亚：}
\BookLine{157.50}{22.14}{}{0}{是啊，尊贵的朋友有时候也该为自己的缺席补}
\BookLine{166.00}{22.14}{}{0}{上点礼物。}
\newpage
% 原书页码：108
\null
\BookNumber{108}
\Vertical{4}{169}{}{6.2}{8}阿\\尔\\刻\\提\\斯\\之\\歌\EndVertical
\BookLine{30.00}{0.00}{\heifont}{0}{看门人：}
\BookLine{38.50}{22.14}{}{0}{阿格莱亚，你这是怎么回事！}
\BookLine{55.50}{0.00}{\heifont}{0}{阿格莱亚：}
\BookLine{64.00}{22.14}{}{0}{{\kaifont （暴躁地）}我知道怎么回事。}
\BookLine{81.00}{7.38}{\kaifont}{0}{两个仆人从宫殿里出来，抬着一张卧榻和一个脚凳。}
\BookLine{98.00}{0.00}{\heifont}{0}{看门人：}
\BookLine{106.50}{22.14}{}{0}{{\kaifont （对仆人）}过来，放这儿，你们知道的。}
\BookLine{123.50}{7.38}{\kaifont}{0}{看门人转过身去扶阿德墨托斯，阿格莱亚指挥了起来。}
\BookLine{140.50}{0.00}{\heifont}{0}{阿格莱亚：}
\BookLine{149.00}{22.14}{}{0}{这儿，放这儿！{\kaifont （照着太阳找到位置）}转过来。}
\BookLine{157.50}{22.14}{}{0}{{\kaifont （从他们手里拿过垫子）}这样，对了。王后喜欢}
\BookLine{166.00}{22.14}{}{0}{坐在这儿。}
\newpage
% 原书页码：109
\null
\BookNumber{109}
\Vertical{129}{11}{}{7}{10}第\\二\\幕\EndVertical
\BookLine{30.00}{7.38}{\kaifont}{0}{阿德墨托斯由人搀扶着走了出来，双手分别搭在看门人}
\BookLine{38.50}{0.00}{\kaifont}{0}{和一个守卫的肩上。}
\BookLine{55.50}{0.00}{\heifont}{0}{阿德墨托斯：}
\BookLine{64.00}{22.14}{}{0}{已经快要正午了。{\kaifont （对看门人说）}你已经守了一}
\BookLine{72.50}{22.14}{}{0}{夜，去休息吧。}
\BookLine{89.50}{0.00}{\heifont}{0}{看门人：}
\BookLine{98.00}{22.14}{}{0}{啊，陛下，像我们这些老年人，睡得都很少。}
\BookLine{106.50}{22.14}{}{0}{轻点，好一点。}
\BookLine{123.50}{0.00}{\heifont}{0}{阿德墨托斯：}
\BookLine{132.00}{22.14}{}{0}{我没看到王后呢，阿格莱亚。}
\BookLine{149.00}{0.00}{\heifont}{0}{阿格莱亚：}
\BookLine{157.50}{22.14}{}{0}{{\kaifont （忙着整理垫子）}陛下，她马上就过来。您放}
\BookLine{166.00}{22.14}{}{0}{心吧。}
\newpage
% 原书页码：110
\null
\BookNumber{110}
\Vertical{4}{169}{}{6.2}{8}阿\\尔\\刻\\提\\斯\\之\\歌\EndVertical
\BookLine{30.00}{7.38}{\kaifont}{0}{阿格莱亚匆忙地走回宫中。}
\BookLine{47.00}{0.00}{\heifont}{0}{看门人：}
\BookLine{55.50}{22.14}{}{0}{陛下，等您舒服地坐好了，我有些消息要告诉}
\BookLine{64.00}{22.14}{}{0}{您。估计您都不敢相信。{\kaifont （斜起眼看看太阳）}是}
\BookLine{72.50}{22.14}{}{0}{的，陛下，今天会非常热。我们已经进入夏至}
\BookLine{81.00}{22.14}{}{0}{时节了。据海员们说，太阳现在很低，在这些}
\BookLine{89.50}{22.14}{}{0}{日子，那个祂离我们非常近。他们就是这么}
\BookLine{98.00}{22.14}{}{0}{说的。}
\BookLine{115.00}{0.00}{\heifont}{0}{阿德墨托斯：}
\BookLine{123.50}{22.14}{}{0}{所以是什么消息？}
\BookLine{140.50}{0.00}{\heifont}{0}{看门人：}
\BookLine{149.00}{22.14}{}{0}{一位尊贵的朋友，陛下，要来见您。一位，好}
\BookLine{157.50}{22.14}{}{0}{朋友。}
\newpage
% 原书页码：111
\null
\BookNumber{111}
\Vertical{129}{11}{}{7}{10}第\\二\\幕\EndVertical
\BookLine{30.00}{0.00}{\heifont}{0}{阿德墨托斯：}
\BookLine{38.50}{22.14}{}{0}{赫拉克勒斯！}
\BookLine{55.50}{0.00}{\heifont}{0}{看门人：}
\BookLine{64.00}{22.14}{}{0}{是的，宙斯之子赫拉克勒斯！}
\BookLine{81.00}{0.00}{\heifont}{0}{阿德墨托斯：}
\BookLine{89.50}{22.14}{}{0}{今天！今天真是个好日子！}
\BookLine{106.50}{0.00}{\heifont}{0}{看门人：}
\BookLine{115.00}{22.14}{}{0}{整个镇上都传开了，他走得也真是够慢的了。}
\BookLine{123.50}{22.14}{}{0}{每到一个村庄，或者农场，人们都献上好酒，}
\BookLine{132.00}{22.14}{}{0}{给他戴上花环。他呢，喝得醉醺醺的，哦，陛}
\BookLine{140.50}{22.14}{}{0}{下，今天可不能让他拥抱您啊。您知道他那力}
\BookLine{149.00}{22.14}{}{0}{气，还不得把您害死。}
\BookLine{166.00}{0.00}{\heifont}{0}{阿德墨托斯：}
\BookLine{174.50}{22.14}{}{0}{王后知道吗？}
\newpage
% 原书页码：112
\null
\BookNumber{112}
\Vertical{4}{169}{}{6.2}{8}阿\\尔\\刻\\提\\斯\\之\\歌\EndVertical
\BookLine{30.00}{0.00}{\heifont}{0}{看门人：}
\BookLine{38.50}{22.14}{}{0}{女仆们应该已经告诉她了。}
\BookLine{55.50}{0.00}{\heifont}{0}{阿德墨托斯：}
\BookLine{64.00}{22.14}{}{0}{他要是换个时间来就好了。}
\BookLine{81.00}{0.00}{\heifont}{0}{看门人：}
\BookLine{89.50}{22.14}{}{0}{赫拉克勒斯想什么时候来就什么时候来。陛下，}
\BookLine{98.00}{22.14}{}{0}{小王子们都激动坏了，伊皮伊尼兹肯定又会让}
\BookLine{106.50}{22.14}{}{0}{他讲一遍杀死复仇女神的狮子和打扫奥革阿斯}
\BookLine{115.00}{22.14}{}{0}{牛棚的故事。}
\BookLine{132.00}{0.00}{\heifont}{0}{阿德墨托斯：}
\BookLine{140.50}{22.14}{}{0}{去把阿格莱亚叫来。}
\BookLine{157.50}{7.38}{\kaifont}{0}{阿格莱亚从宫殿里走出来。}
\newpage
% 原书页码：113
\null
\BookNumber{113}
\Vertical{129}{11}{}{7}{10}第\\二\\幕\EndVertical
\BookLine{30.00}{0.00}{\heifont}{0}{看门人：}
\BookLine{38.50}{22.14}{}{0}{她来了。}
\BookLine{55.50}{0.00}{\heifont}{0}{阿格莱亚：}
\BookLine{64.00}{22.14}{}{0}{陛下，您有什么吩咐？}
\BookLine{81.00}{7.38}{\kaifont}{0}{阿格莱亚背着国王打手势让看门人先离开。看门人从左}
\BookLine{89.50}{0.00}{\kaifont}{0}{侧下场。}
\BookLine{106.50}{0.00}{\heifont}{0}{阿德墨托斯：}
\BookLine{115.00}{22.14}{}{0}{阿格莱亚，来。你听说了吗？}
\BookLine{132.00}{7.38}{\kaifont}{0}{看门人刚要下场，很紧张地看了阿格莱亚一眼，听到她}
\BookLine{140.50}{0.00}{\kaifont}{0}{的回答后才算放心。}
\BookLine{157.50}{0.00}{\heifont}{0}{阿格莱亚：}
\BookLine{166.00}{22.14}{}{0}{听说了，宙斯之子，要来看望我们。}
\newpage
% 原书页码：114
\null
\BookNumber{114}
\Vertical{4}{169}{}{6.2}{8}阿\\尔\\刻\\提\\斯\\之\\歌\EndVertical
\BookLine{30.00}{0.00}{\heifont}{0}{阿德墨托斯：}
\BookLine{38.50}{22.14}{}{0}{{\kaifont （淡然地）}阿格莱亚，这是我最后的日子了。我}
\BookLine{47.00}{22.14}{}{0}{知道，我感觉到了。}
\BookLine{64.00}{0.00}{\heifont}{0}{阿格莱亚：}
\BookLine{72.50}{22.14}{}{0}{陛下，这种事情没人会知道的。您还是少说几}
\BookLine{81.00}{22.14}{}{0}{句吧。您需要保存体力。}
\BookLine{98.00}{0.00}{\heifont}{0}{阿德墨托斯：}
\BookLine{106.50}{22.14}{}{0}{太阳落山之前我就会死去。如果可行的话，我}
\BookLine{115.00}{22.14}{}{0}{希望你能办到，拖上个一两天，不要让赫拉克}
\BookLine{123.50}{22.14}{}{0}{勒斯得知我死亡的消息。你就说我们有个仆人}
\BookLine{132.00}{22.14}{}{0}{死了，就像什么都没发生过一样照常生活。王}
\BookLine{140.50}{22.14}{}{0}{后会告诉你该怎么做。你知道赫拉克勒斯一直}
\BookLine{149.00}{22.14}{}{0}{以来的生活，总是担着沉重的工作，一个接一}
\BookLine{157.50}{22.14}{}{0}{个。他高高兴兴地来，我不想让他看到一片哀}
\BookLine{166.00}{22.14}{}{0}{伤。尤其是赫拉克勒斯，总是出生入死，他很}
\BookLine{174.50}{22.14}{}{0}{讨厌葬礼啊哀哭啊这些事情。}
\newpage
% 原书页码：115
\null
\BookNumber{115}
\Vertical{129}{11}{}{7}{10}第\\二\\幕\EndVertical
\BookLine{30.00}{0.00}{\heifont}{0}{阿格莱亚：}
\BookLine{38.50}{22.14}{}{0}{{\kaifont （又整理起垫子）}陛下，我不傻，您说的这些我}
\BookLine{47.00}{22.14}{}{0}{都明白。您好好躺下，闭上眼睛休息，阳光会}
\BookLine{55.50}{22.14}{}{0}{带给您力量的。}
\BookLine{72.50}{0.00}{\heifont}{0}{阿德墨托斯：}
\BookLine{81.00}{22.14}{}{0}{再给我一天的力气就好，好来欢迎赫拉克勒斯。}
\BookLine{89.50}{22.14}{}{0}{你去看看王后在哪里呢。}
\BookLine{106.50}{0.00}{\heifont}{0}{阿格莱亚：}
\BookLine{115.00}{22.14}{}{0}{她来了。}
\BookLine{132.00}{7.38}{\kaifont}{0}{阿格莱亚从左侧下场。阿尔刻提斯从宫殿里走出来。她}
\BookLine{140.50}{0.00}{\kaifont}{0}{站在最高一级的台阶上，穿着第一幕的那条裙子。阿德墨托}
\BookLine{149.00}{0.00}{\kaifont}{0}{斯回头看到她，没有说话。}
\BookLine{166.00}{0.00}{\heifont}{0}{阿德墨托斯：}
\BookLine{174.50}{22.14}{}{0}{这是咱们结婚时你穿的衣服！}
\newpage
% 原书页码：116
\null
\BookNumber{116}
\Vertical{4}{169}{}{6.2}{8}阿\\尔\\刻\\提\\斯\\之\\歌\EndVertical
\BookLine{30.00}{7.38}{\kaifont}{0}{阿尔刻提斯微笑着，手指放在唇边示意他不要说话，并}
\BookLine{38.50}{0.00}{\kaifont}{0}{走到他身边，轻声说着。}
\BookLine{55.50}{0.00}{\heifont}{0}{阿尔刻提斯：}
\BookLine{64.00}{22.14}{}{0}{妻子们都留着它等着葬礼的时候穿，我穿上它}
\BookLine{72.50}{22.14}{}{0}{是为了好好活着。}
\BookLine{89.50}{0.00}{\heifont}{0}{阿德墨托斯：}
\BookLine{98.00}{22.14}{}{0}{亲爱的，坐这儿来。我刚还跟阿格莱亚说……}
\BookLine{115.00}{7.38}{\kaifont}{0}{阿尔刻提斯坐在阿德墨托斯身边的凳子上，把手放在他}
\BookLine{123.50}{0.00}{\kaifont}{0}{的手边。}
\BookLine{140.50}{0.00}{\heifont}{0}{阿尔刻提斯：}
\BookLine{149.00}{22.14}{}{0}{是的，我知道，等赫拉克勒斯来了……}
\BookLine{166.00}{0.00}{\heifont}{0}{阿德墨托斯：}
\BookLine{174.50}{22.14}{}{0}{他马上就到了。}
\newpage
% 原书页码：117
\null
\BookNumber{117}
\Vertical{129}{11}{}{7}{10}第\\二\\幕\EndVertical
\BookLine{30.00}{0.00}{\heifont}{0}{阿尔刻提斯：}
\BookLine{38.50}{22.14}{}{0}{一定是在路上耽搁了，只有危险的时候他才跑}
\BookLine{47.00}{22.14}{}{0}{得飞快，说不定他要今晚才能到，或者要等明}
\BookLine{55.50}{22.14}{}{0}{天了。不过我们不着急。{\kaifont （起身走到舞台中央，}}
\BookLine{64.00}{22.14}{}{0}{{\kaifont 站在那里倾听着）}我听到山谷里传来的欢呼声}
\BookLine{72.50}{22.14}{}{0}{了。{\kaifont （走回到凳子旁边，重复着）}我们不着急。}
\BookLine{81.00}{22.14}{}{0}{{\kaifont （站在舞台中央，望着太阳仿佛要寻求帮助）}这}
\BookLine{89.50}{22.14}{}{0}{样的阳光能够治愈人，夏至日的阳光，你感觉}
\BookLine{98.00}{22.14}{}{0}{到了吗？}
\BookLine{115.00}{0.00}{\heifont}{0}{阿德墨托斯：}
\BookLine{123.50}{22.14}{}{0}{这阳光能治愈人，但只能治愈别人。我已经不}
\BookLine{132.00}{22.14}{}{0}{需要这些了，不再需要什么希望。我的一生很}
\BookLine{140.50}{22.14}{}{0}{短暂，但短短一个钟头也足以成为完美的时光。}
\BookLine{149.00}{22.14}{}{0}{我得到过了，最完美的时光。一个拥有过幸福}
\BookLine{157.50}{22.14}{}{0}{的人，就不会成为时间的囚徒。来，在最后这}
\BookLine{166.00}{22.14}{}{0}{一天，让我们把还没说完的话都好好说完。}
\end{document}
